\documentclass[10pt]{article}

\usepackage[utf8]{inputenc}

\usepackage[T1]{fontenc}

\usepackage{amsmath}

\usepackage{amsfonts}

\usepackage{amssymb}

\usepackage[version=4]{mhchem}

\usepackage{stmaryrd}

\usepackage{fixltx2e}

\usepackage{graphicx}

\usepackage[export]{adjustbox}

\graphicspath{ {./images/} }



\begin{document}

Справедливы тождества



\[

\left.\begin{array}{l}

z\left[J_{v}(z) N_{v}^{\prime}(z)-N_{v}(z) J_{v}^{\prime}(z)\right]=\frac{2}{\pi}, \\

z\left[H_{v}^{(1)}(z) H_{v}^{(2)^{\prime}}(z)-H_{v}^{(2)}(z) H_{v}^{(1)^{\prime}}(z)\right]=-\frac{4 i}{\pi}

\end{array}\right\}

\]



и соотнотения



\[

\begin{aligned}

& J_{v}(z)=\frac{1}{2}\left[H_{v}^{(1)}(z)+H_{v}^{(2)}(z)\right], \\

& H_{v}(z)=\frac{1}{2 i}\left[H_{v}^{(1)}(z)-H_{v}^{(2)}(z)\right] .

\end{aligned}

\]



Для действительных $z=x$ и $v$ функции Ганкеля являются комплексно сопряженными решениями уравнения (1). При этом функции $J_{v}(z)$ дают действительную часть, а функции $N_{v}(x)$ - мнимую часть функций Ганкеля.



Ц. Ф. 1-го, 2-го и 3-го рода $Z_{v}$ удовлетворяют рекуррентным формулам



\[

\begin{aligned}

& Z_{v-1}(z)+Z_{v+1}(z)=\frac{2 v}{z} Z_{v}(z), \\

& Z_{v-1}(z)-Z_{v+1}(z)=2 Z_{v}^{\prime}(z) .

\end{aligned}

\]



Каждая пара функций



\[

J_{v}(z), J_{-v}(z) ; J_{v}(z), Y_{v}(z) ; H_{v}^{(1)}(z), H_{v}^{(2)}(z)

\]



образует (при нецелом $v$ ) фундаментальную систему решіений уравнения (1).



М о д и ф ци р о в а н ны м и Ц. ф. наз. Ц. ф. мнимого аргумента



\[

I_{v}(z)=\left\{\begin{array}{l}

e^{-i v \pi / 2} J_{v}\left(e^{i \pi / 2} z\right),-\pi<\arg z \leqslant \pi / 2, \\

e^{-3 i v \pi / 2} J_{v}\left(e^{-3 i \pi / 2} z\right), \pi / 2<\arg z \leqslant \pi,

\end{array}\right.

\]



и Макдональда функции:



\[

K_{v}(z)=\frac{1}{2} i \pi e^{i \pi v / 2} H_{v}^{(1)}\left(e^{i \pi / 2} z\right)=

\][



=-\textbackslash frac\{1\}\{2\} i \textbackslash pi e\^{}\{-i \textbackslash pi v / 2\} H\_\{v\}\textsuperscript{\{(2)\}\textbackslash left(e}\{-i \textbackslash pi / 2\} z\textbackslash right)=\textbackslash frac\{1\}\{2\} i \textbackslash pi e\^{}\{-i v \textbackslash pi / 2\} H\_\{v\}\textsuperscript{\{(1)\}\textbackslash left(e}\{i \textbackslash pi / 2\} z\textbackslash right) \textbackslash text \{. \}

]



Эти функции являются решениями дифференциального уравнения



\[

z^{2} \frac{d^{2} Z}{d z^{2}}+z \frac{d Z}{d z}-\left(z^{2}+v^{2}\right) Z=0

\]



и удовлетворяют рекуррентным формулам



\[

\begin{gathered}

I_{v-1}(z)-I_{v+1}(z)=\frac{2 v}{z} I_{v}(z), \\

K_{v-1}(z)-K_{v+1}(z)=-\frac{2 v}{z} K_{v}(z), \\

K_{-v}(z)=K_{v}(z) .

\end{gathered}

\]



Цилиндрические функции целых и полуцелых порядков. Если $v=n$ - целое число, то $J_{n}(z)$ можно ошределить с помощью фо р м у лы Я к о и-А н г е р а



\[

\exp \left\{\frac{z}{2}\left(t-\frac{1}{t}\right)\right\}=\sum_{n=-\infty}^{+\infty} t^{n} J_{n}(z)

\]



или



\[

\begin{gathered}

\exp \{ \pm i z \sin \theta\}=J_{0}(z)+ \\

+2 \sum_{n=1}^{\infty}\left\{J_{2 n}(z) \cos 2 n \theta \pm i J_{2 n+1}(z) \sin (2 n+1) \theta\right\} .

\end{gathered}

\]



Справедливы равенства



\[

\begin{gathered}

J_{n}(z)=\sum_{m=0}^{\infty} \frac{(-1)^{m} z^{n+2 m}}{2^{n+2 m} m !(m+n) !}, \\

J_{-n}(z)=(-1)^{n} J_{n}(z) .

\end{gathered}

\]



Функция $J_{n}(z)$ есть целая трансцендентная функция аргумента $z ;$ для алгебрапческого $z=a, a \neq 0, J_{n}(z)$ есть трансцендентное число и $J_{m}(a) \neq J_{n}(a)$ при $m \neq n$. В качестве второго линейно независимого с $J_{n}(z)$ решения уравнения (1) обычно берут функцию



\[

\begin{gathered}

N_{n}(z)=\lim _{v \rightarrow n} \frac{1}{\sin v \pi}\left[J_{v}(z) \cos v \pi-J_{-v}(z)\right]= \\

=\frac{1}{\pi} \lim _{v \rightarrow n}\left[\frac{\partial J_{v}(z)}{\partial v}-(-1)^{n} \frac{\partial J_{-v}(z)}{\partial z}\right]=\frac{2}{\pi}\left(c+\ln \frac{z}{2}\right) J_{n}(z)- \\

-\frac{1}{\pi}\left(\frac{z}{2}\right)^{n} \sum_{m=0}^{\infty}\left\{\frac{(-1)^{m}}{m !(m+n) !}\left(\frac{z}{2}\right)^{2 m} \times\right. \\

\left.\times\left(\sum_{l=1}^{m} \frac{1}{l}+\sum_{l=1}^{m+n} \frac{1}{l}\right)\right\}- \\

-\frac{1}{\pi}\left(\frac{z}{2}\right)^{-n} \sum_{m=0}^{n-1} \frac{n-m-1}{m !}\left(\frac{z}{2}\right)^{2 m}, n=0,1,2, \ldots, \text { (8) }

\end{gathered}

\]



где $c=0,577215 \ldots-$ постоянная Эйлера. Если в одной из конечных сумм верхний индекс суммирования меньше нижнего, то соответствующая сумма получает значение 0. Справедливо равенство



\[

Y_{-n}(z)=(-1)^{n} Y_{n}(z)

\]



Ц. Ф. тогда и только тогда превращаются в элементарные функции, когда индекс $v$ принимает значение $v=n+1 / 2, n=0,1,2, \ldots$ (с ф е р и ч е с к и е ф у н кци и Бе с с е ля, или Ц. ф. по луцелого пор я д к а). Справедливы формулы $(n=0,1,2, \ldots)$ :



\[

J_{n+1 / 2}(z)=(-1)^{n} \sqrt{\frac{2}{\pi}} z^{n+1 / 2}\left(\frac{1}{z} \frac{d}{d z}\right)^{n} \frac{\sin z}{z},

\]



в частности $J_{1 / 2}(z)=\sqrt{\frac{\overline{2}}{\pi z}} \sin z$



\[

J_{-n-1 / 2}(z)=\sqrt{\frac{2}{\pi}} z^{n+1 / 2}\left(\frac{1}{z} \frac{d}{d z}\right)^{n} \frac{\cos z}{z},

\]



в частности $J_{-1 / 2}(z)=\sqrt{\frac{2}{\pi z}} \cos z$;



\[

\begin{gathered}

H_{-n-1 / 2}^{(1)}(z)=i(-1)^{n} H_{n+1 / 2}^{(1)}(z), \\

H_{-n-1 / 2}^{(2)}(z)=-i(-1)^{n} H_{n+1 / 2}^{(2)}(z), \\

N_{n+1 / 2}(z)=(-1)^{n+1} J_{-n-1 / 2}(z) ; \\

N_{-n-1 / 2}(z)=(-1)^{n} J_{n+1 / 2}(z) ; \\

K_{n+1 / 2}(z)=K_{-n-1 / 2}(z)=(-1)^{n} \sqrt{\frac{\pi}{2 z}} z^{n+1}\left(\frac{d}{z d z}\right)^{n} \frac{e-z}{z} .

\end{gathered}

\]



Интегральные представления цилиндрических функций. Для $v=n=0,1,2, \ldots$ имеется и н т е г р а л ь-



\begin{center}

\includegraphics[max width=\textwidth]{2023_07_09_bd9adb925a58d3c32a2eg-1}

\end{center}



\[

J_{n}(z)=\frac{1}{\pi} \int_{0}^{\pi} \cos (z \sin \varphi-n \varphi) d \varphi

\]



II



\[

N_{n}(z)=\frac{1}{\pi} \int_{0}^{\pi} \sin (z \sin \varphi-n \varphi) d \varphi-

\][



-\textbackslash frac\{2\}\{\textbackslash pi\} e\^{}\{-i n \textbackslash pi / 2\} \textbackslash int\_\{0\}\^{}\{\textbackslash infty\} e\^{}\{-z \textbackslash operatorname\{sh\} t\} \textbackslash operatorname\{cth\} n\textbackslash left(t+\textbackslash frac\{1\}\{2\} i \textbackslash pi\textbackslash right) d t, R(z)>0 .

]



Для $R\left(0+\frac{1}{2}\right)>0$ и $R(z)>0 \quad$ имеется ин т е гральное представление Пуассона



\[

\begin{gathered}

J_{v}(z)=\frac{2\left(\frac{1}{2} z\right)^{v}}{\sqrt{\pi} \Gamma\left(v+\frac{1}{2}\right)} \int_{0}^{\pi / 2} \cos (z \sin \varphi) \cos ^{2 v} \varphi d \varphi, \\

I_{v}(z)=\frac{\left(\frac{1}{2} z\right)^{v}}{\sqrt{ } \bar{\pi} \Gamma\left(v+\frac{1}{2}\right)} \int_{-1}^{1} e^{-z t}\left(1-t^{2}\right)^{v-1 / 2} d t

\end{gathered}

\]





\end{document}